\chapter{Rezultate}
\label{cap:evaluare}

\section{Setup și metrici}

Sistemul este evaluat pe un set de aur de 98 de cazuri (35 naștere, 30 căsătorie, 33 locuire), dintre care 16 de tip \emph{refuz controlat} (în afara domeniului); proveniența și validarea sa sunt discutate în Secțiunea~\ref{sec:limitari}. Toate rulările folosesc aceiași parametri ($k=4$, o singură reîncercare, temperatura 0) și aceeași infrastructură (GPU 2$\times$Tesla T4). Comparația cantitativă principală este un \textbf{studiu comparativ} a șase configurații: modelul 3B la trei niveluri de cuantizare (FP16, Q8, Q4) și două scări mai mari la Q4 (7B și 14B), plus un model RO-specializat din altă familie (RoLlama3-8B), rulat prin pipeline-ul complet. Codul sursă, setul de aur și artefactele de evaluare sunt disponibile public la \url{https://github.com/cosmin-glod/Romanian-Legal-Document-Assistent}, pentru reproductibilitate.

Se calculează cinci metrici: \emph{refusal accuracy} (\texttt{is\_refusal} vs \texttt{should\_refuse}), \emph{breadcrumb recall@k}, \emph{lexical keyword recall} (potrivire exactă de sub-șir a cuvintelor-cheie așteptate), \emph{rată de validare} (cazuri care trec toate regulile R1-R6) și \emph{latency}. Recall-ul de breadcrumb și cel de cuvinte-cheie sunt \emph{proxy-uri}: al doilea este o limită inferioară a fidelității (nu recunoaște sinonime), triangulată mai jos printr-un eșantion evaluat manual.

\section{Rezultate principale}

Tabelul~\ref{tab:sweep} sintetizează rezultatele celor șase configurații: modelul 3B la trei niveluri de cuantizare (FP16, Q8, Q4), două scări mai mari la Q4 (7B și 14B) și modelul RO-specializat RoLlama3-8B; fidelitatea răspunsurilor este evaluată separat de un model judecător de 14B.

\begin{table}[tb]
    \centering\footnotesize
    \begin{tabular}{l c c c c c c}
        \toprule
        \textbf{Config} & \textbf{Refusal} & \textbf{Recall} & \textbf{Lexical} & \textbf{Validare} & \textbf{Fidelit.} & \textbf{Lat.\,(s)} \\
        \midrule
        \texttt{3B-Q4}   & 62 & 95 & 41 & 53 & 76 & \textbf{13} \\
        \texttt{3B-Q8}   & 74 & 95 & 53 & 68 & 86 & 16 \\
        \texttt{3B-FP16} & 79 & 95 & 55 & 68 & 83 & 27 \\
        \texttt{7B-Q4}   & 69 & 95 & 57 & 68 & 92 & 28 \\
        \texttt{14B-Q4}  & 66 & 95 & 47 & \textbf{71} & \textbf{100} & 37 \\
        \midrule
        \texttt{RoLlama-8B} & \textbf{82} & 95 & \textbf{71} & 56 & 79 & 48 \\
        \bottomrule
    \end{tabular}
    \caption{Rezultate pe setul de aur: cuantizare, scară și model RO-specializat.}
    \label{tab:sweep}
\end{table}

Se remarcă trei aspecte principale. Întâi, cuantizarea la 4 biți are un cost real: \texttt{3B-Q4} este cea mai slabă configurație de 3B la refuz, acoperire lexicală și validare, în timp ce \texttt{3B-Q8} se apropie de \texttt{3B-FP16} la o fracțiune din memorie. Apoi, precizia poate compensa scara: \texttt{3B-Q8} egalează sau depășește la refusal accuracy modelele mai mari \texttt{7B-Q4} și \texttt{14B-Q4}, deci un model mic, bine cuantizat, este competitiv cu unul mare slab cuantizat. În al treilea rând, fidelitatea (judecată de modelul de 14B) crește monoton cu scara, în timp ce recall-ul rămâne constant pe toate configurațiile, ceea ce confirmă că factorul determinant este generarea, nu regăsirea. Judecătorul nu a marcat aproape niciun răspuns drept \emph{incorect}, deci valorile de fidelitate sunt o limită superioară optimistă, triangulată de eșantionul evaluat manual mai jos.

Schimbarea familiei de model aduce o observație distinctă. RoLlama3-8B, specializat pe limba română, reduce drastic \emph{refuzul nejustificat} (3 cazuri, față de 33 pe \texttt{3B-Q4}), ceea ce confirmă că specializarea lingvistică acționează direct asupra modului de eșec dominant (Secțiunea~\ref{sec:nejustificat}). Compromisul constă într-o deplasare către eșecul opus, \emph{sub-refuzul} (modelul răspunde la 15 din 16 întrebări din afara domeniului), și într-o disciplină mai slabă de citare și ancorare, surprinsă de stratul de verificare prin încălcări ale R2, R4 și R6. Cu alte cuvinte, un model specializat pe română deplasează profilul de eroare dinspre supra-refuz către sub-refuz și erori de ancorare, adică tocmai zona în care verificarea ieșirii aduce valoare.

Pentru a calibra lexical keyword recall (41\% pe 3B), 15 răspunsuri non-refuz au fost evaluate manual: 27\% corecte, 47\% parțiale, 27\% incorecte. Metrica lexicală subevaluează deci calitatea, însă fidelitatea pe 3B rămâne modestă, eroarea recurentă fiind nu inventarea, ci \emph{confuzia de rol} (proprietar $\leftrightarrow$ chiriaș, drepturi $\leftrightarrow$ obligații), pe care verificarea nu o identifică.

\section{Stratul de verificare: erori nesemnalate prinse}

Rolul principal al verificării este detectarea răspunsurilor plauzibile, dar nesusținute de surse, adică a erorilor nesemnalate descrise în \cite{gheorghe2026robust}. Regula R4 (documente inventate) domină încălcările pe 3B (26 de apariții) și se prăbușește odată cu scara (2 pe 7B, 0 pe 14B), în timp ce R5 (refuz însoțit de documente inexistente) apare doar pe 7B (8 apariții). Astfel, creșterea scării schimbă clasa de eșec, surprinsă de aceeași infrastructură. Precizia regulii R4 a fost auditată manual: din 10 cazuri semnalate, 9 reprezintă probleme reale ($\approx$ 90\% precizie).

O \emph{ablație} pe configurația \texttt{Qwen2.5-3B-Q4} măsoară direct valoarea stratului: în absența verificării, \textbf{66 din 98} de răspunsuri (67\%) ar fi fost servite cu defect, în majoritate din cauza lipsei citațiilor (60 de răspunsuri neverificabile). Reîncercarea \textbf{repară 48} dintre ele și \textbf{reține 18} sub forma unui refuz controlat, astfel încât niciun răspuns defect nu ajunge la cetățean. Structural, citarea (R2) este o precondiție pentru R4: în prima încercare aproape niciun răspuns nu conține citații, astfel încât documentele inventate (R4) devin vizibile abia după ce reîncercarea adaugă citațiile.

\section{Modul de eșec dominant: refuzul nejustificat}
\label{sec:nejustificat}

Cea mai frecventă eroare nu este inventarea de informație, ci \emph{refuzul nejustificat}: modelul răspunde ,,Nu am suficiente informații" la o întrebare al cărei răspuns este, de fapt, prezent în surse. Pe \texttt{3B-Q4}, \textbf{33 din 82} de întrebări din domeniu (40\%) primesc un refuz nejustificat, dintre care 31 aveau răspunsul regăsit în top-$k$; pe \texttt{7B-Q4}, \textbf{28 din 82} (34\%). Refuzul nejustificat depășește inventarea de documente și rămâne modul de eșec dominant. Scara nu îl reduce semnificativ (\texttt{14B-Q4} păstrează 30 de cazuri), în schimb o cuantizare de calitate da (\texttt{3B-Q8} coboară la 20). Un prompt reformulat recuperează doar 13 din 33 de cazuri, deci eșecul ține de calibrarea generării, nu de regăsire.

Prin construcție, stratul de verificare nu poate detecta acest mod de eșec: un refuz trece de R1-R6 fără să conțină vreo citație sau document de verificat. Această limitare a motivat regula \textbf{R7}, care semnalează un refuz suspect atunci când similaritatea top-1 depășește un prag. Calibrat pe setul de aur, pragul 0.64 separă refuzurile corecte ($\leq$0.63) de cele nejustificate (până la 0.80); R7 detectează \textbf{24 din 33} de cazuri (73\%) cu \textbf{zero fals pozitive}, fiind mai precis decât reformularea promptului și independent de generare. În practică, regula este folosită pentru \emph{semnalarea} cazurilor în vederea revizuirii, nu pentru recuperarea automată.

Pentru ilustrare, mai jos sunt date un răspuns corect și un eșec tipic, ambele din rularea pe \texttt{7B-Q4}:

\begin{center}
\setlength{\fboxsep}{7pt}%
\fcolorbox{codeframe}{codebg}{%
\begin{minipage}{0.92\textwidth}
\small\raggedright
\textbf{Răspuns corect (ancorat).} \emph{Întrebare:} ,,Ce documente îmi trebuie pentru înregistrarea nașterii copilului, dacă sunt căsătorit cu partenera?" \emph{Răspuns:} ,,Actele de identitate ale părinților; Certificatul medical constatator al nașterii\ldots", cu citarea surselor S1 și S3.\\[5pt]
\textbf{Refuz nejustificat.} \emph{Întrebare:} ,,Cine poate primi indemnizația pentru creșterea copilului?" \emph{Răspuns:} ,,Nu am suficiente informații în sursele disponibile pentru a răspunde.", deși răspunsul se afla în fragmentul regăsit cu cea mai mare similaritate.
\end{minipage}}
\end{center}

\section{Evaluarea pre-completării formularelor}

Stratul de pre-completare a fost evaluat pe 14 conversații (două per formular, 180 de câmpuri), cu o singură rundă de extracție deterministă. Acuratețea globală este \textbf{99\%} (179/180): șase din șapte formulare ating 100\%, singurul sub acest prag fiind contractul de închiriere (97\%). Un pas determinist suplimentar (expresii regulate pentru date în format ISO sau românesc și deducerea sexului din „băiat"/„fetiță") a crescut acuratețea de la 94\% la 99\% și rata de completare dintr-o singură rundă de la 57\% la 93\%. Singurul câmp ratat provine dintr-o confuzie de rol (adresa proprietarului în câmpul imobilului); deoarece niciun câmp nu este inventat (lipsă $\to$ null $\to$ completare manuală), fiecare formular conduce la un PDF complet.

